\documentclass[11pt]{article}
\usepackage[T1]{fontenc}
\usepackage{lmodern}
\usepackage[letterpaper,top=.93in,bottom=.94in,left=.98in,right=.98in]{geometry}
\usepackage{amsmath,amssymb,amsthm}
\usepackage{microtype}
\usepackage[unicode,hidelinks]{hyperref}
\hypersetup{pdftitle={Five interlacing rows of ordered set partitions at threshold four},pdfauthor={Bowen Zuo}}
\newtheorem{theorem}{Theorem}
\newcommand{\e}{\varepsilon}
\newcommand{\E}{\mathbb E}
\newcommand{\ip}{\mathrm i}
\title{Five interlacing rows of ordered set partitions at threshold four}
\author{Bowen Zuo\\East China University of Science and Technology}
\date{November 2026}
\begin{document}
\maketitle

\begin{abstract}
Let $O_n(t)$ enumerate ordered set partitions of $[n]$ by the number of
blocks of size at least four. For every $n\ge0$, the rows $O_n$ through
$O_{n+4}$ have simple negative roots and satisfy strict pairwise
interlacing, with constants interpreted by empty root lists. Five is the
largest such window uniform in $n$. We prove the assertion for $n\ge288$
by following the four zeros of a meromorphic denominator in a fixed
disk. A cancellation at an actual zero of the old row gives the sign
for shift four, even when the leading oscillation of the two rows is
nearly in phase opposition. Independent exact integer computations
certify the remaining $1152$ comparisons.
\end{abstract}

\section{The ordered model and the root order}
An ordered set partition is a list of disjoint nonempty blocks with
union $[n]$. The empty list contributes $O_0=1$. Every block of size at
least four contributes one factor $t$. Removing the first block gives
\begin{equation}\label{eq:recurrence}
 O_n(t)=\sum_{a=1}^n \binom{n}{a} t^{[a\ge4]} O_{n-a}(t).
\end{equation}
Write $P_j(z)=\sum_{a=0}^j z^a/a!$, $A=P_3-1$ and $B=e^z-P_3$.
Thus the exponential generating function in the number of labels is
\begin{equation}\label{eq:egf}
 \sum_{n\ge0}O_n(t)\frac{z^n}{n!}=\frac{1}{1-A(z)-tB(z)}.
\end{equation}
For $0\le k\le\lfloor n/4\rfloor$, $k$ blocks of size four and the
remaining singletons give a positive coefficient of $t^k$.
There are no higher coefficients. For example,
$O_4=74+t$, $O_5=530+11t$, and $O_8=521640+24125t+70t^2$.

\begin{theorem}\label{thm:main}
For every integer $n\ge0$ and $h\in\{1,2,3,4\}$, all zeros of $O_n$
and $O_{n+h}$ are simple and negative. Put $d=\lfloor n/4\rfloor$
and list their respective roots as $a_i$ and $b_i$ in increasing
order. If both rows have degree $d$, then
\[
 a_1<b_1<a_2<b_2<\cdots<a_d<b_d<0 .
\]
If $O_{n+h}$ has degree $d+1$, then
\[
 b_1<a_1<b_2<a_2<\cdots<b_d<a_d<b_{d+1}<0 .
\]
For $d=0$ these statements mean, respectively, two empty root lists
or the single inequality $b_1<0$. No six-row window has this
pairwise interlacing property uniformly in its starting rank.
\end{theorem}

The ordered large-block question appears in the Symmetric Functions
Catalog \cite{catalog}. Fang's Theorem~1 \cite{fang} proves the
eventual five-row theorem for the unordered threshold-four model,
whose generating function is $\exp(A+tB)$. Fang centers the fourth
shift at an actual old-row zero immediately after Eq.~(8) in
\cite[\S5.1]{fang} and matches absolute zero counts in \cite[\S6]{fang}.
The analogous steps below use residues and an explicit endpoint.
Fang's coefficient order concerns the unordered rows
\cite[\S8]{fang}. Multiplying their $t^k$ coefficients by $k!$
orders the marked large blocks, giving $\exp(A)/(1-tB)$.
Equation~\eqref{eq:egf} orders every block. We use
\eqref{eq:recurrence} for the finite certificates.

\section{Four poles in a fixed contour}
For $t<0$ put $q=-t/(1-t)$, $\delta=1-q$ and $\e=(24\delta)^{1/4}$.
Equation~\eqref{eq:egf} gives
\begin{equation}\label{eq:G}
 \frac{O_n(t)}{n!}=\delta[z^n]G_q(z)^{-1},\qquad
 G_q(z)=q e^z+2\delta-P_3(z).
\end{equation}
At any zero $z$ of $G_q$, one has $G'_q(z)=z^3/6-2\delta$.
On the real axis the derivative changes from negative to positive
at the unique solution of $q=P_2(x)e^{-x}$, since the derivative
of $P_2(x)e^{-x}$ is $-x^2 e^{-x}/2$. The minimum is $-2J(x)$,
where
\[
 J(x)=P_2(x)e^{-x}-1+x^3/12,qquad
 J'(x)=x^2(1/4-e^{-x}/2).
\]
There is a unique positive zero $\alpha$ of $J$, with
$947/1000<\alpha<948/1000$ by rational Taylor bounds.
Set $q_c=1-\alpha^3/12>116/125$ and
$\e_c=(2\alpha^3)^{1/4}<1143/1000$. For $q_c<q<1$ the minimum is
positive. At $q_c$ the only real zero is $\alpha$, where
$G'_{q_c}(\alpha)=0$ and $G''_{q_c}(\alpha)=\alpha^2/2>0$.

Let $U$ have density $4(1-u)^3$ on $[0,1]$ and put $H(z)=\E e^{Uz}$.
Then
\begin{equation}\label{eq:Gsplit}
 G_q(z)=qz^4H(z)/24+\delta(2-P_3(z)),\qquad
 \E U^a=24a!/(a+4)! .
\end{equation}
For $|\Im z|\le3$, the real part of
$e^{-\ip\Im z/2}H(z)$ is positive. Indeed
$|(U-1/2)\Im z|\le3/2<\pi/2$. On $|z|=3$ we also have
$\Re H(z)>3/10$. If $x=\Re z\le0$, the bound
$\cos v\ge1-v^2/2$ gives
\[
 \Re H(z)\ge\E e^{-3U}-\tfrac92\E U^2
 =\tfrac{24}{81}(2+e^{-3})-\tfrac3{10}>\tfrac3{10}.
\]
For $0\le x\le3$ use $1+x/5-(9-x^2)\E[U^2e^{xU}]/2$.
The expectations at $x=1,2,3$ are strictly below $1/8,1/5,3/10$,
respectively, by the positive Taylor bounds in the arithmetic
certificate. On $[0,1]$, $[1,2]$, and $[2,3]$, replacing the
expectation by the corresponding upper bound and using
$x\ge0,1,2$ gives the stated real part bound.
Since $|2-P_3(z)|\le13$ on this circle, Rouch\'e's theorem counts
four zeros with multiplicity inside it for $q_c\le q\le1$.
It also gives
\begin{equation}\label{eq:circle}
 \min_{|z|=3}|G_q(z)|>
 \tfrac{81}{24}\tfrac{3}{10}q-13(1-q)\ge \tfrac1{300}.
\end{equation}
The lower bound exceeds $1$ when $\e\le1/10$.
The imaginary axis has
$\Re G_q(\ip y)=\delta+y^2/2-q(1-\cos y)>0$ for $q<1$.
As $q\uparrow1$ the four scaled roots tend to the fourth roots
of $-1$, as proved below. The open right and left half-planes
each contain two zeros counted with multiplicity for $q_c<q<1$.
No zero crosses the imaginary axis or the contour at $q_c$.
There the right pair is the double zero $\alpha$ and the left
pair is nonreal and conjugate. On $(q_c,1)$ the real axis is
zero-free and each conjugate pair is simple. Write $\zeta=\e w$
and $\eta=\e v$ for the upper right and upper left zeros. They
are continuous and locally analytic.

\subsection{Geometry of the scaled roots}\label{sec:boxes}
In scaled coordinates the equation is
\[
 F_{\e}(w)=2-e^{\e w}+w^4H(\e w)=0.
\]
For $\omega^4=-1$ one root lies in
$|w-(\omega-3\omega^2\e/10)|<r\e^2$
for $(H_0,r)=(1/32,1/2)$ or $(1/10,3/5)$ and $0<\e\le H_0$.
On the boundary set $u=3/10+rH_0$ and $W=1+uH_0$.
In $F_{\e}=1+w^4+\e(w^5/5-w)+O(\e^2)$ the constant and
first-order terms at the shifted center cancel. Compare with
$4\omega^3(w-\omega+3\omega^2\e/10)$. The quartic
linearization remainder, change in the first-order term, and
exponential remainder, divided by $\e^2$, are bounded by
\begin{equation}\label{eq:disk-small}
 (6+4uH_0+(uH_0)^2)u^2+u(W^4+1)
    +\frac{W^2+W^6/15}{2(1-H_0W)}<4r.
\end{equation}
Here the expectation identities in \eqref{eq:Gsplit} and
$e^x\le(1-x)^{-1}$ for $0\le x<1$ bound the exponential
remainders. The disks are disjoint and map inside $|z|<3$.
The squared moduli of the upper right and left centers are
$1\mp 3\sqrt2\e/10+9\e^2/100$, respectively.
For the upper right center the real part exceeds the imaginary
part by $3\e/10$. The bounds $r\e^2$ on displacement preserve
these strict modulus and angular inequalities.

For $1/32\le\e\le1143/1000$, the data file
\texttt{dominance.json} covers the entire closed range by 17
closed rational bins sharing endpoints. Each bin gives an upper left disk
outside $|w|=1$ and a right disk inside $|w|=1$ with
$\Re w>\Im w$. The right disk is in the upper half-plane except
for the last bin, where it is centered on the real axis and
contains two zeros with multiplicity. On the oscillatory side
these are exactly the two right conjugates. All disks map
inside $|z|<3$.

For clarity, the disk tests have the following exact meaning.
For a rational center $c_0$ and radius $r_0$, at the parameter
midpoint $e_0$ compare $F_\e$ on the circle with
$F_{e_0}^{(\ell)}(c_0)(w-c_0)^\ell/\ell!$, where $\ell=1$
for a single zero and $\ell=2$ for the last right disk.
The jet through order two is summed through degree 28 in
Gaussian rational arithmetic. For monomial index $a>0$ the
coefficient is
\[
 \frac{-e_0^a+\mathbf1_{a\ge4}\,24e_0^{a-4}}{a!};
\]
the term with $e_0^{a-4}$ is used only for $a\ge4$.
The jet tail is bounded by twice its first absolute majorant
at $|w|\le2$, since consecutive ratios are below $1/10$.
On $|w|\le W$ a parameter derivative bound is
$W e^{\e W}(1+W^4/5)$. The next $w$-derivative
bound follows from Leibniz's rule on $w^4H(\e w)$.
Expectations of $U^b e^{\e W U}$ are bounded by positive
Taylor sums through degree 20 with a geometric tail.
The verifier compares a lower norm bound for the leading jet
with the sum of the other jet, disk remainder and parameter
variation bounds. This proves the strict Rouch\'e inequality
on the whole cell. Together these tests give
\begin{equation}\label{eq:geometry}
 R:=|\zeta|<\e<|\eta|,\qquad
 0<\theta:=\arg\zeta<\pi/4 ,\qquad q_c<q<1.
\end{equation}

\section{Signs at actual old-row zeros}
Let $\phi=\arg\eta$, $\beta=\arg G'_q(\zeta)$,
$\gamma=\arg G'_q(\eta)$, $\lambda=R/|\eta|$ and
$A_0=|G'_q(\zeta)|/|G'_q(\eta)|$.
Use the principal $\beta\in(0,\pi)$, continuous because
$\Im G'_q(\zeta)=R^3\sin(3\theta)/6>0$.
For $\psi=(n+1)\theta+\beta$ and $\chi=(n+1)\phi+\gamma$,
the residue theorem gives
\begin{align}
 S_n&:=\tfrac12 |G'_q(\zeta)|R^{n+1}[z^n]G_q(z)^{-1}
 =-\cos\psi-A_0\lambda^{n+1}\cos\chi+E_n,\label{eq:residue}\\
 |E_n|&\le \frac{3}{2\min_{|z|=3}|G_q(z)|}
        |G'_q(\zeta)|(R/3)^{n+1}.\label{eq:E}
\end{align}
The normalization preserves the sign of $O_n(t)$. At an old
zero, put $c=T+E$ with
$T\ge A_0\lambda^{n+1}$ and $E\ge\max(|E_n|,|E_{n+h}|)$.
Thus $|\cos\psi|\le c$. It is enough to prove
\begin{equation}\label{eq:tests}
 c<2/5,\quad c^2+L_h^2<1,\quad
 |S_{n+h}-\sin\psi\sin(h\theta)|\le L_h\sin(h\theta),\quad
 L_h\ge0 .
\end{equation}
Indeed $\sin(h\theta)>0$ and $|\sin\psi|>L_h$.
For $h=1,2,3$ the absolute error in \eqref{eq:tests} is at most
$2T+2E$ using \eqref{eq:residue} at the old root.
For $h=4$ add the identity $S_n=0$ to \eqref{eq:residue} at
$n+4$. As $\cos(\chi+3\pi)=-\cos\chi$, using the same
$\chi$ at both rows gives the sharper bound
\begin{equation}\label{eq:cancel}
 c|1+\cos4\theta|+
 T(1-\lambda^4+|4\phi-3\pi|)+2E.
\end{equation}
This is the residue version of the centered shift-four comparison
discussed above. Bounding these two secondary terms individually
would lose the small angular factor.

\subsection{The range \texorpdfstring{$0<\e\le1/10$}{0 < epsilon <= 1/10}}
Fix $n\ge288$ and $n\e\ge3$. For each upper scaled pole
write its ratio to $\omega$ as $1+d_0$. The disks give
$|d_0+3\omega\e/10|\le r\e^2$, $|d_0|\le u\e$.
In a cell $[l,H_1]$, put $u=3/10+rH_1$ and
$B_1=r+u^2/(2(1-uH_1))$. The principal logarithm differs
from $-3\omega\e/10$ by at most $B_1\e^2$.
With $L=.2121320-B_1H_1$ and $U_1=.2121321+B_1H_1$,
these logarithmic estimates imply
\[
 4L\e\le\pi-4\theta\le4U_1\e,\quad
 |4\phi-3\pi|\le4U_1\e,\quad
 2L\e\le-\log\lambda\le2U_1\e .
\]
The two decimal constants enclose $3/(10\sqrt2)$. The sine
bounds are $\sin(4\theta)\ge\e S$, where
$S=4L-(4U_1)^3H_1^2/6$, and
$\sin(h\theta)\ge .7-3U_1H_1$ for $1\le h\le3$.
The latter uses the Lipschitz bound on sine and
$\sin(h\pi/4)\ge\sqrt2/2>.7$.

At a pole $G'_q(\e w)=\e^3(w^3-\e/2)/6$.
Put $M=(1+uH_1)^3+H_1/2$ and $m=(1-uH_1)^3-H_1/2$.
Since $n\e\ge\max(3,288l)$, we may use
\[
 T=\frac{M}{m}
    \left(\sum_{a=0}^{20}\frac{(2L\max(3,288l))^a}{a!}\right)^{-1}.
\]
The contour bound for this range allows
$|E_n|, |E_{n+h}|\le P\e$ with $P=MH_1^2/4$,
so choose $E=PH_1$. For the short shifts use
$L_h=(2T+2E)/(.7-3U_1H_1)$. For the fourth shift,
$1-\lambda^4\le -4\log\lambda$ and
$1+\cos4\theta\le (4U_1\e)^2/2$ give
\begin{equation}\label{eq:small-four}
 L_4=\frac{12U_1T+c(4U_1)^2H_1/2+2P}{S}.
\end{equation}
The exact verifier checks positive denominators and both strict
tests in \eqref{eq:tests} on 256 equal cells closed at both
ends. Their boundary at $1/32$ permits using the sharper disk
radius on the first 80 cells.

\subsection{The compact range and the critical cell}
Use the 16 bins from Section~\ref{sec:boxes} that meet
$[1/10,\e_c]$; the remaining bin is contained in the small
range. In a bin $[l,H_1]$ let $(x,y),(X,Y)$ be the centers
and $r,s$ the radii of the right and left disks, respectively.
Let $R_l,R_u,L_l,L_u$ enclose the two scaled moduli by
outward rational square-root rounding. Then take
\begin{align*}
 \Lambda&=R_u/L_l<1, &\lambda_0&=R_l/L_u>0,\\
 d_-&=l^3(L_l^3-H_1/2)/6, &d_+&=H_1^3(R_u^3+H_1/2)/6,\\
 T&=\frac{R_u^3+H_1/2}{L_l^3-H_1/2}\Lambda^{288},&
 E&=450d_+(H_1R_u/3)^{288}.
\end{align*}
In the ratio for $T$ the factor $\e^3/6$ cancels at the same
parameter. The quantities $d_-,d_+$ bound the left and right
physical residue derivatives from below and above.
Equation~\eqref{eq:circle} gives the factor 450.

For a simple upper right disk write $x_l=x-r$, $y_l=y-r$.
The short-shift sine bound is $y_l/R_u$. Also
\begin{align*}
 \sin4\theta&\ge
 4x_ly_l(x-y-2r)(x+y-2r)/R_u^4,\\
 \frac{1+\cos4\theta}{\sin4\theta}
 &\le\frac{(x-y+2r)(x+y+2r)}{2x_ly_l},\\
 |4\phi-3\pi|&\le 4(|X+Y|+2s)/(Y-X-2s).
\end{align*}
The first two use the Cartesian sine and cosine identities,
and the third uses the arctangent of the deviation from
the upper left diagonal. These inequalities and
$1-\lambda^4\le1-\lambda_0^4$ give the short and fourth
relative errors in \eqref{eq:tests}.

In the final right disk write $\zeta=x_0+\ip y_0$ with
$0<x_0,y_0<6/5$ at an oscillatory parameter.
The imaginary part of $G_q(\zeta)=0$ yields
\begin{equation}\label{eq:critical}
 \Re G'_q(\zeta)
 =(y_0\cot y_0-1)(1+x_0+x_0^2/2-y_0^2/6)+y_0^2/3 .
\end{equation}
The factor in parentheses after $y_0\cot y_0-1$ lies
between 0 and 3. Use
$\sin y_0-y_0\cos y_0=\int_0^{y_0} t\sin t\,dt\le y_0^3/3$
and $\sin y_0\ge3y_0/4$ to obtain
$|\Re G'_q(\zeta)|\le5y_0^2/3$.
The imaginary part is $y_0(3x_0^2-y_0^2)/6$ and has
absolute value at most $y_0$. In particular
$|G'_q(\zeta)|\le5y_0$. For $1\le h\le3$,
$\sin(h\theta)\ge\sin\theta$ gives
$|G'_q(\zeta)|/\sin(h\theta)\le 5H_1R_u$.
For $h=4$ use
$\sin4\theta=4\sin\theta\cos\theta\cos2\theta$,
$\cos\theta>1/2$ and $\cos2\theta\ge(x_0-y_0)/R$.
The ratio is then at most $5(H_1R_u)^2/(2l(x-y-2r))$.
If $D_1$ is the larger of these bounds, take
\[
 L_h=2D_1\left(\Lambda^{288}/d_-+
                 450(H_1R_u/3)^{288}\right)
 \qquad(1\le h\le4).
\]
This follows by bounding the two secondary and contour
terms directly, retaining the factor $|G'_q(\zeta)|$.
The inequality is uniform up to the critical parameter
from the open side. The verifier checks \eqref{eq:tests}
with exact rational arithmetic. Across all the small cells
and compact bins, the largest squared-test values display
as $0.987357440$ and $0.064608115$, respectively.

\section{Counting to the degree}
The right poles tend to $\alpha$ as $q\downarrow q_c$.
Equation~\eqref{eq:critical} and the imaginary derivative
formula give $\Re G'_q(\zeta)=O(y_0^2)$ and
$\Im G'_q(\zeta)/y_0\to\alpha^2/2>0$.
Consequently $\psi\to\pi/2$.
Stop at $q_*=1-(3/n)^4/24$, so $\e=3/n\le1/96$ and
$n\e\ge3$ on the intervening interval.
For $\omega=e^{\ip\pi/4}$, the derivative ratio at the endpoint is
\[
 \frac{G'_q(\e w)}{\e^3\omega^3/6}
 =(1+d_0)^3-\e/(2\omega^3).
\]
With $H_*=3/288$, $r=1/2$, $u=3/10+rH_*$ and
$D_*=(3+3uH_*+u^2H_*^2)u+1/2$, this ratio is within
$D_*\e$ of 1. Hence
$|\beta-3\pi/4|\le D_*\e/(1-D_*H_*)$.
Let $B_*$ be the log remainder constant above at $H_*$ and
$b_*=.2121321$. The endpoint phase loss from
$(n/4+1)\pi$ is at most
\[
 3b_*+H_*\left(b_*+(3+H_*)B_*+D_* /(1-D_*H_*)\right)
 <.671<\pi/4 .
\]
The displayed rational expression is strictly below $671/1000$;
$n/4$ in this phase estimate is real division.

By continuity, the phase crosses in order the $d=\lfloor n/4\rfloor$
strips $[i\pi+\pi/4,(i+1)\pi-\pi/4]$ for $i=1,\dots,d$.
For each strip take its first upper endpoint after the preceding
strip and the last lower endpoint before that upper endpoint.
This gives disjoint parameter intervals in which the phase stays
within the corresponding strips. At opposite endpoints
$-\cos\psi$ has opposite signs and magnitude $\cos(\pi/4)>2/5$.
The bound $|S_n+\cos\psi|\le c<2/5$ from the same estimates
holds at both endpoints. Thus there is a zero in each interval.
The known degree makes these $d$ roots simple and exhaustive.

At the $i$th such zero in increasing $q$, the shifted row has
sign $(-1)^i$ by \eqref{eq:tests}. At $q=0$ it is positive.
There is one shifted zero before the first old zero and one
between each pair of consecutive old zeros. When the degree
rises, positivity of the leading coefficient gives the sign
$(-1)^{d+1}$ near $q=1$, producing the additional last root.
Again degree counts show simplicity and exhaustiveness.
As $t=-q/(1-q)$ decreases with $q$, the two root orders become
exactly those in Theorem~\ref{thm:main}. This proves its assertion
for all old ranks $n\ge288$.

\section{Finite certificates and maximality}
For $0\le n<288$, there are $1152$ comparisons, with upper rank
at most 291. The ancillary finite generator implements
\eqref{eq:recurrence} over the integers and locates candidate
rational root brackets. The independent verifier computes the
coefficient of $t^k$ in the original row using
\[
 [t^k]O_n(t)=n![z^n]B(z)^k(1-A(z))^{-k-1}.
\]
This expression follows by recording the $k$ ordered large
blocks and the $k+1$ lists of small blocks between and around them.
A direct enumeration by first-block subsets checks the
initial rows through rank 7 independently. The verifier checks
strictly negative endpoints, disjointness of each row's
brackets, and opposite signs by integer Horner evaluation with
positive cleared denominators. The bracket count equals the
degree, so each interval has one simple root, and the intervals
account for every root. Compare their endpoints across each
pair to obtain the specified strict order.
There are $10512$ isolating intervals, 426 equal-positive-degree
comparisons, 710 comparisons with positive old degree increased
by one, 6 with two constants and 10 from a constant to a linear
row. These facts and the analytic tail prove all-sizes
interlacing. Every five consecutive rows form such a window.
For $n=4m+3$ the degrees of $O_n$ and $O_{n+5}$ differ by two,
which rules out a uniform six-row window.

\paragraph{Reproducibility.}
The ancillary directories contain the rational pole tests and
the independent integer verifier for the finite range.
The companion Lean~4 project proves Theorem~\ref{thm:main} for
the original rows. In the namespace
\texttt{OrderedPartitions.FourPole}, the declaration
\texttt{original\_all} covers all old ranks and shifts one through
four, \texttt{original\_five\_rows} covers all pairs in five
consecutive rows, and \texttt{six\_rows\_fail} proves sharpness.
The file \texttt{FourPoleAllSizes.lean} joins the two proofs at
old rank 288. The reproducibility materials identify the fixed Lean
and Mathlib versions and the complete source archive.

\paragraph{Acknowledgment.}
ChatGPT assisted with the analytic arguments, independent
mathematical review, exact certificate implementation, Lean
formalization, and manuscript preparation.

\begin{thebibliography}{9}
\bibitem{fang}
C. Fang, \emph{Eventual interlacing of large-block set-partition
polynomials}, 26 November 2026,
\href{https://hexagonmath.org/2611.01129v1}{hexagon:2611.01129v1}.
\bibitem{catalog}
P. Alexandersson, \emph{Real-rooted and stable polynomials},
The Symmetric Functions Catalog, set partitions and trees,
\href{https://www.symmetricfunctions.com/realRootedOpenProblems.htm\#set-partitions-and-trees}{catalog entry}.
\end{thebibliography}
\end{document}
